%mac
\begin{dang}{PHƯƠNG PHÁP GIẢI CÁC BÀI TOÁN ĐỘNG LỰC HỌC}
\begin{itemize}
\item[1.] Chọn vật khảo sát chuyển động. Biểu diễn các lực tác dụng lên vật, trong đó làm rõ phương, chiều và điểm đặt của từng lực.
\item[2.] Chọn hai trục vuông góc Ox và Oy; trong đó trục Ox cùng hướng với chuyển động của vật hay cùng hướng với lực kéo khi vật đứng yên. Phân tích các lực theo hai trục này. Áp dụng định luật 2 Newton theo hai trục toạ độ Ox và Oy.
\begin{align*}
\begin{cases}
&Ox: F_x=F_{1x}+F_{2x}+...=m.a_x\quad (1)\\
&Oy:  F_y=F_{1y}+F_{2y}+...=m.a_y\quad (2)
\end{cases}
\end{align*}
\item[3.] Giải hệ phương trình $(1)$ và $(2)$ để tìm gia tốc hay tìm lực, tuỳ từng bài toán.
\end{itemize}
\end{dang}